\documentclass[11pt]{article}

\usepackage[margin=1in]{geometry}
\usepackage{amsmath,amssymb,amsthm}

\newtheorem*{claim}{Claim}
\newtheorem*{lemma}{Lemma}

\begin{document}

\title{Load Comp.}
\author{}
\date{}
\maketitle

\begin{claim}
Consider the problem
\[
        P||\sum_j p_j C_j(\alpha) 
\]

And for any instance $I$ consider a list schedule $S^W$ that maximizes the objective and a 
list schedule $S^O$ that minimizes the objective. Let $W_i$ and $O_i, i \in [m]$
be their final machine loads. Then
\[
        \min_i W_i \leq O_k
        \qquad \forall k  \in [m].
\]
\end{claim}

\begin{proof}
We use again the form derived by SChweigelsohn
Writing $L_i$ for the final load of machine $i$ in a list schedule S of instance I, 
summing over all machines
gives
\[
        \sum_j p_j C_j
        =
        \frac12\sum_i L_i(\sigma(I))^2
        +
        \frac12\sum_j p_j^2 .
\]
The second term is independent of the list schedule. Hence, among list
schedules, minimizing $\sum_j p_j C_j$ is equivalent to minimizing
$\sum_i L_i(\sigma(I))^2$, and maximizing $\sum_j p_j C_j$ is equivalent to maximizing
$\sum_i L_i(\sigma(I))^2$.

It remains to prove the following exchange property of list schedules.

\begin{lemma}[Conditional on the exchange statement]
Let $W=(W_1,\dots,W_m)$ be the final load vector of a list schedule $\sigma^{W}$ maximizing
the objective for instance $I$, and let $O=(O_1,\dots,O_m)$ be the final load vector of a list
schedule $S^O$ minimizing the objective for instance $I$. After sorting both vectors in nonincreasing
order,
\[
        W_1\geq W_2\geq \cdots \geq W_m,
        \qquad
        O_1\geq O_2\geq \cdots \geq O_m,
\]
assume that every violated partial-sum cut of the minimum-square schedule
\(S^O\) admits a non-overshooting exchange. Then
\[
        \sum_{i=1}^r W_i
        \geq
        \sum_{i=1}^r O_i
        \qquad
        \text{for } r=1,\dots,m
\]
and equality holds for $r=m$ because both schedules are of the same instance.
\end{lemma}

\begin{proof}
We make the exchange argument explicit. Let \(X=(X_1,\dots,X_m)\) be the
final load vector of some list schedule. The elementary operation to keep in
mind is the transfer of one job of size \(p\) from machine \(a\) to machine
\(b\). If before the transfer the two relevant loads are \(X_a\) and \(X_b\),
then after the transfer they are
\[
        X_a-p,
        \qquad
        X_b+p.
\]
Hence the change in the square-load objective is
\[
\begin{aligned}
\Delta
&=
(X_a-p)^2+(X_b+p)^2-X_a^2-X_b^2  \\
&=
2p(X_b-X_a)+2p^2  \\
&=
-2p(X_a-X_b-p).
\end{aligned}
\]
Thus, if \(X_a-X_b>p\), moving the job from the larger-load machine \(a\) to
the smaller-load machine \(b\) strictly decreases \(\sum_i X_i^2\). Conversely,
moving the same job from the smaller-load machine to the larger-load machine
strictly increases \(\sum_i X_i^2\), since in that case
\[
        (X_a+p)^2+(X_b-p)^2-X_a^2-X_b^2
        =
        2p(X_a-X_b)+2p^2
        >
        0.
\]

In a list schedule, an adjacent interchange can be analyzed as a finite sequence
of such elementary transfers between the machines whose assigned jobs change.
The list rule assigns the earlier job to a machine whose load at that time is
minimal. Therefore the sign of the change in the square-load objective is
controlled by the direction and size of these transfers: transfers from smaller
final loads to larger final loads increase \(\sum_i X_i^2\), while transfers
from larger final loads to smaller final loads decrease \(\sum_i X_i^2\)
whenever the moved job does not overshoot the load gap. The inequalities above
give the exact sign condition for each elementary transfer.

Now suppose that the claimed majorization inequality were false. Then there is
some \(r<m\) such that
\[
        \sum_{i=1}^r W_i
        <
        \sum_{i=1}^r O_i .
\]
Both schedules process the same total load, say \(P\), so
\[
        \sum_{i=r+1}^m W_i
        =
        P-\sum_{i=1}^r W_i
        >
        P-\sum_{i=1}^r O_i
        =
        \sum_{i=r+1}^m O_i .
\]
Thus, compared with \(W\), the vector \(O\) has an excess among its \(r\)
largest machines and the same amount of deficit among its \(m-r\) smallest
machines. Write this excess as
\[
        E
        =
        \sum_{i=1}^r O_i-\sum_{i=1}^r W_i
        >
        0.
\]
At this point one has to be careful: the displayed inequalities do not, by
themselves, imply that there is an indivisible job whose size does not overshoot
the relevant load gap. Instead we get the following useful dichotomy.

Let a feasible adjacent interchange move a net amount \(x>0\) from a machine
\(a\le r\) in the larger-load block to a machine \(b>r\) in the smaller-load
block. For a single transfer \(x=p\). For a swap of a job of size \(p\) with a
job of size \(q\), \(x=p-q\). Put
\[
        d=O_a-O_b.
\]
After the exchange, the two relevant loads are
\[
        O_a' = O_a-x,
        \qquad
        O_b' = O_b+x,
\]
and all other loads are unchanged. The partial-sum excess across the cut
\(\{1,\dots,r\}\) is reduced by exactly \(x\):
\[
\begin{aligned}
E'
&=
\left(\sum_{i=1}^r O_i-x\right)
-
\sum_{i=1}^r W_i  \\
&=
E-x
<
E.
\end{aligned}
\]
At the same time, the square-load objective changes by
\[
\begin{aligned}
\sum_i (O_i')^2-\sum_i O_i^2
&=
(O_a-x)^2+(O_b+x)^2-O_a^2-O_b^2 \\
&=
-2x(d-x).
\end{aligned}
\]
Now there are two cases.

If
\[
        0<x<d,
\]
then
\[
        -2x(d-x)<0.
\]
Thus the exchange strictly decreases \(\sum_i X_i^2\), contradicting the
choice of \(O\) as a minimum-square list schedule.

If no such non-overshooting exchange exists, then every feasible exchange
across this cut has \(x\ge d\). A strict overshoot \(x>d\) increases the
square-load objective by
\[
        2x(x-d).
\]
Equivalently, because
\[
        \sum_j p_j C_j(S)
        =
        \frac12\sum_i L_i(S)^2+\frac12\sum_j p_j^2,
\]
the objective increases by
\[
        x(x-d).
\]
This is the direction relevant for an upper bound on the gap from best to
worst. Write
\[
        F(S)=\sum_j p_j C_j(S).
\]
Fix any sequence of adjacent interchanges
\[
        S^0=S^O,\quad S^1,\quad \dots,\quad S^T=S^W
\]
that transforms the list order of \(S^O\) into the list order of \(S^W\), and
write
\[
        \Delta_\ell=F(S^\ell)-F(S^{\ell-1}).
\]
Then the desired objective gap telescopes as
\[
        F(S^W)-F(S^O)
        =
        \sum_{\ell=1}^T \Delta_\ell
        \leq
        \sum_{\ell=1}^T \max\{\Delta_\ell,0\}.
\]
Thus only objective-increasing exchanges can contribute to the upper bound on
``worst minus best''. For the \(\ell\)-th exchange of the type above, let
\(x_\ell>0\) be the net amount transferred from the larger-load side to the
smaller-load side, and let \(d_\ell\) be the corresponding load gap before the
exchange. Then
\[
        \Delta_\ell=x_\ell(x_\ell-d_\ell).
\]
Consequently \(\Delta_\ell>0\) only in the strict overshoot case
\(x_\ell>d_\ell\), and we get the upper bound
\[
\begin{aligned}
F(S^W)-F(S^O)
&\le
\sum_{\ell:\,x_\ell>d_\ell} x_\ell(x_\ell-d_\ell)  \\
&\le
\sum_{\ell:\,x_\ell>d_\ell} x_\ell^2 .
\end{aligned}
\]
Thus, if a violated cut cannot be repaired by a non-overshooting exchange, then
the contribution of that obstruction to the worst-best gap is controlled by the
overshoot penalties \(x_\ell(x_\ell-d_\ell)\), not by a lower bound on the gap.

There is one useful statement that can be proved directly. For a sorted load
vector \(X\), write
\[
        P_r(X)=\sum_{i=1}^r X_i .
\]
Suppose a single transfer changes \(X\) by moving \(x>0\) from a machine \(a\)
that is among the \(r\) largest loads of \(X\) to a machine \(b\) that is not
among the \(r\) largest loads of \(X\). Put \(d=X_a-X_b\). If \(x\ge d\), then
after the transfer the new load of \(b\) is
\[
        X_b+x\ge X_a .
\]
Thus, in the sorted vector after the transfer, the load \(X_b+x\) can replace
the old top-\(r\) load \(X_a\) by a load that is at least as large. The other
\(r-1\) old top loads are unchanged. Therefore
\[
        P_r(X')\ge P_r(X).
\]
Consequently, any transfer that actually decreases the sorted top-\(r\)
partial sum must satisfy
\[
        0<x<d.
\]
So along any transformation path from a schedule with
\(P_r(X)>P_r(W)\) to \(S^W\), the first step that lowers \(P_r\) is necessarily
a non-overshooting transfer.

This proves the non-overshooting property for any step that decreases the
violated partial sum, but it does not yet prove that such a step is available
directly from \(S^O\). That distinction matters: previous steps along the path
could have increased the objective before the first partial-sum-decreasing step
occurs. Hence, to finish the majorization proof, one still needs the stronger
exchange statement that every violated partial-sum cut of the minimum-square
schedule \(S^O\) itself admits a non-overshooting exchange.

If that stronger exchange statement is established, then no such \(r\) exists:
the exchange would strictly decrease \(\sum_i X_i^2\), contradicting the choice
of \(O\) as a minimum-square list schedule. Therefore
\[
        \sum_{i=1}^r W_i
        \geq
        \sum_{i=1}^r O_i
        \qquad
        r=1,\dots,m-1.
\]
For \(r=m\), both sides are equal to the total processing time of the instance.
\end{proof}

Assuming the exchange statement, apply the lemma. Since $W\succ O$ and the
total loads are equal, we have
\[
        \sum_{i=1}^{m-1} W_i
        \geq
        \sum_{i=1}^{m-1} O_i .
\]
Subtracting from the common total processing time gives
\[
        W_m \leq O_m .
\]
Because the vectors were sorted in nonincreasing order,
\[
        W_m=\min_i W_i,
        \qquad
        O_m=\min_i O_i.
\]
Therefore
\[
        \min_i W_i
        \leq
        \min_i O_i.
\]
Finally, every machine load in the optimal schedule is at least the minimum
load in that schedule, so for every machine $k$,
\[
        \min_i W_i
        \leq
        \min_i O_i
        \leq
        O_k.
\]
Thus the claim follows once the exchange statement used in the lemma is proved.
\end{proof}

\end{document}
